\subsection{Numerical Operator Architecture and Precedence Graphs}

Numerical operators are not magic syntax pasted on top of raw memory. In Python, an operator is an expression-level request sent to runtime objects. The objects involved decide whether the operation is supported and what kind of result object should be produced.

For example:

\begin{verbatim}
result = 40 + 2

print(f"result: {result}")
print(f"type:   {type(result)}")
\end{verbatim}

Possible output:

\begin{verbatim}
result: 42
type:   <class 'int'>
\end{verbatim}

The expression \texttt{40 + 2} produces a new integer object. The name \texttt{result} is then bound to that object.

\subsubsection{Standard Arithmetic: Operator Dispatch and Result Object Creation}

The usual arithmetic operators are:

\begin{verbatim}
+    addition
-    subtraction
*    multiplication
/    true division
//   floor division
%    remainder
**   exponentiation
\end{verbatim}

A first example:

\begin{verbatim}
a = 42
b = 8

print(f"a + b:  {a + b}")
print(f"a - b:  {a - b}")
print(f"a * b:  {a * b}")
print(f"a / b:  {a / b}")
print(f"a // b: {a // b}")
print(f"a % b:  {a % b}")
print(f"a ** 2: {a ** 2}")

print()
print(f"type(a + b):  {type(a + b)}")
print(f"type(a / b):  {type(a / b)}")
print(f"type(a // b): {type(a // b)}")
\end{verbatim}

Possible output:

\begin{verbatim}
a + b:  50
a - b:  34
a * b:  336
a / b:  5.25
a // b: 5
a % b:  2
a ** 2: 1764

type(a + b):  <class 'int'>
type(a / b):  <class 'float'>
type(a // b): <class 'int'>
\end{verbatim}

Operators are connected to special methods on objects.

\begin{verbatim}
number = 42

print(f"type(number): {type(number)}")

print()
print("selected operator methods:")
print(f"{'__add__':<15} {'__add__' in dir(number)}")
print(f"{'__sub__':<15} {'__sub__' in dir(number)}")
print(f"{'__mul__':<15} {'__mul__' in dir(number)}")
print(f"{'__truediv__':<15} {'__truediv__' in dir(number)}")
print(f"{'__floordiv__':<15} {'__floordiv__' in dir(number)}")
print(f"{'__mod__':<15} {'__mod__' in dir(number)}")
print(f"{'__pow__':<15} {'__pow__' in dir(number)}")

print()
print(f"number.__add__(8): {number.__add__(8)}")
print(f"number + 8:        {number + 8}")
\end{verbatim}

Possible output:

\begin{verbatim}
type(number): <class 'int'>

selected operator methods:
__add__        True
__sub__        True
__mul__        True
__truediv__    True
__floordiv__   True
__mod__        True
__pow__        True

number.__add__(8): 50
number + 8:        50
\end{verbatim}

In normal code, we do not call \texttt{\_\_add\_\_()} directly. We write \texttt{number + 8}. The method inspection only shows the object-protocol machinery behind the operator.

Mixed numeric types may produce a wider result type:

\begin{verbatim}
int_value = 42
float_value = 0.5
complex_value = 2j

result1 = int_value + float_value
result2 = int_value + complex_value
result3 = float_value + complex_value

print(f"result1: {result1}, type={type(result1)}")
print(f"result2: {result2}, type={type(result2)}")
print(f"result3: {result3}, type={type(result3)}")
\end{verbatim}

Possible output:

\begin{verbatim}
result1: 42.5, type=<class 'float'>
result2: (42+2j), type=<class 'complex'>
result3: (0.5+2j), type=<class 'complex'>
\end{verbatim}

The name on the left is not typed. It simply receives the result object produced by the expression.

\subsubsection{Division Divergence: Floating-Point Division (\texttt{/}) vs. Floor Division (\texttt{//})}

Python has two different division operators.

\texttt{/} performs true division and produces a floating-point result for ordinary integer operands:

\begin{verbatim}
a = 10
b = 2

result = a / b

print(f"result: {result}")
print(f"type:   {type(result)}")
\end{verbatim}

Possible output:

\begin{verbatim}
result: 5.0
type:   <class 'float'>
\end{verbatim}

Even though the mathematical result is whole, the result object is a \texttt{float}.

\texttt{//} performs floor division:

\begin{verbatim}
a = 11
b = 2

result = a // b

print(f"result: {result}")
print(f"type:   {type(result)}")
\end{verbatim}

Possible output:

\begin{verbatim}
result: 5
type:   <class 'int'>
\end{verbatim}

Floor division means the result is rounded down toward negative infinity, not merely chopped toward zero.

\begin{verbatim}
print(f"11 // 2:   {11 // 2}")
print(f"-11 // 2:  {-11 // 2}")
print(f"11 // -2:  {11 // -2}")
print(f"-11 // -2: {-11 // -2}")
\end{verbatim}

Possible output:

\begin{verbatim}
11 // 2:   5
-11 // 2:  -6
11 // -2:  -6
-11 // -2: 5
\end{verbatim}

This is a common trap for C programmers. In many C contexts, integer division truncates toward zero. Python's \texttt{//} floors.

With floating-point operands, floor division still floors, but the result is a float:

\begin{verbatim}
a = 11.0
b = 2.0

result = a // b

print(f"result: {result}")
print(f"type:   {type(result)}")
\end{verbatim}

Possible output:

\begin{verbatim}
result: 5.0
type:   <class 'float'>
\end{verbatim}

The practical rule is:

\begin{quote}
Use \texttt{/} when you want true division. Use \texttt{//} when you want floor division. Do not treat \texttt{//} as C-style integer truncation.
\end{quote}

\subsubsection{Modular Math (\texttt{\%}) and Exponentiation (\texttt{**}) Mechanics}

The modulo operator \texttt{\%} returns the remainder connected to floor division.

\begin{verbatim}
a = 42
b = 8

quotient = a // b
remainder = a % b

print(f"a:         {a}")
print(f"b:         {b}")
print(f"quotient:  {quotient}")
print(f"remainder: {remainder}")

print()
print(f"check: {quotient} * {b} + {remainder} = {quotient * b + remainder}")
\end{verbatim}

Possible output:

\begin{verbatim}
a:         42
b:         8
quotient:  5
remainder: 2

check: 5 * 8 + 2 = 42
\end{verbatim}

Python keeps this relationship:

\begin{verbatim}
a == (a // b) * b + (a % b)
\end{verbatim}

This also applies with negative values:

\begin{verbatim}
pairs = [
    (11, 2),
    (-11, 2),
    (11, -2),
    (-11, -2),
]

for a, b in pairs:
    q = a // b
    r = a % b

    print(f"a={a:>4}, b={b:>3} | "
          f"a//b={q:>4}, a%b={r:>3} | "
          f"check={q * b + r:>4}")
\end{verbatim}

Possible output:

\begin{verbatim}
a=  11, b=  2 | a//b=   5, a%b=  1 | check=  11
a= -11, b=  2 | a//b=  -6, a%b=  1 | check= -11
a=  11, b= -2 | a//b=  -6, a%b= -1 | check=  11
a= -11, b= -2 | a//b=   5, a%b= -1 | check= -11
\end{verbatim}

The function \texttt{divmod()} gives both results at once:

\begin{verbatim}
a = 42
b = 8

q, r = divmod(a, b)

print(f"quotient:  {q}")
print(f"remainder: {r}")
print(f"type(q):   {type(q)}")
print(f"type(r):   {type(r)}")
\end{verbatim}

Possible output:

\begin{verbatim}
quotient:  5
remainder: 2
type(q):   <class 'int'>
type(r):   <class 'int'>
\end{verbatim}

Exponentiation uses \texttt{**}:

\begin{verbatim}
base = 2
power = 10

result = base ** power

print(f"result: {result}")
print(f"type:   {type(result)}")
\end{verbatim}

Possible output:

\begin{verbatim}
result: 1024
type:   <class 'int'>
\end{verbatim}

With integers, large exponentiation can produce arbitrarily large exact integers:

\begin{verbatim}
big = 2 ** 100

print(f"big: {big}")
print(f"type: {type(big)}")
print(f"bit length: {big.bit_length()}")
\end{verbatim}

Possible output:

\begin{verbatim}
big: 1267650600228229401496703205376
type: <class 'int'>
bit length: 101
\end{verbatim}

A negative exponent produces a floating-point result:

\begin{verbatim}
result = 2 ** -3

print(f"result: {result}")
print(f"type:   {type(result)}")
\end{verbatim}

Possible output:

\begin{verbatim}
result: 0.125
type:   <class 'float'>
\end{verbatim}

So \texttt{**} is still operator dispatch. The result type depends on the operation and operands.

\subsubsection{Unary Operators and Precedence Traps: \texttt{-x}, \texttt{+x}, and \texttt{-2 ** 2}}

Python has unary numeric operators:

\begin{verbatim}
+x
-x
\end{verbatim}

Unary plus usually returns the numeric value unchanged:

\begin{verbatim}
x = 42

print(f"+x: {+x}")
print(f"type(+x): {type(+x)}")
\end{verbatim}

Possible output:

\begin{verbatim}
+x: 42
type(+x): <class 'int'>
\end{verbatim}

Unary minus negates the value:

\begin{verbatim}
x = 42

print(f"-x: {-x}")
print(f"type(-x): {type(-x)}")
\end{verbatim}

Possible output:

\begin{verbatim}
-x: -42
type(-x): <class 'int'>
\end{verbatim}

These operators are also connected to special methods:

\begin{verbatim}
x = 42

print(f"{'__pos__':<12} {'__pos__' in dir(x)}")
print(f"{'__neg__':<12} {'__neg__' in dir(x)}")

print()
print(f"x.__pos__(): {x.__pos__()}")
print(f"x.__neg__(): {x.__neg__()}")
\end{verbatim}

Possible output:

\begin{verbatim}
__pos__     True
__neg__     True

x.__pos__(): 42
x.__neg__(): -42
\end{verbatim}

The famous precedence trap is:

\begin{verbatim}
result = -2 ** 2

print(result)
\end{verbatim}

Possible output:

\begin{verbatim}
-4
\end{verbatim}

This surprises many students because they read it as:

\begin{verbatim}
(-2) ** 2
\end{verbatim}

But Python reads it as:

\begin{verbatim}
-(2 ** 2)
\end{verbatim}

A direct comparison:

\begin{verbatim}
a = -2 ** 2
b = (-2) ** 2
c = -(2 ** 2)

print(f"-2 ** 2:    {a}")
print(f"(-2) ** 2:  {b}")
print(f"-(2 ** 2):  {c}")
\end{verbatim}

Possible output:

\begin{verbatim}
-2 ** 2:    -4
(-2) ** 2:  4
-(2 ** 2):  -4
\end{verbatim}

Exponentiation binds more tightly than unary minus on its left side.

There is a second detail: exponentiation binds less tightly than unary operators on its right side.

\begin{verbatim}
result = 2 ** -2

print(f"2 ** -2: {result}")
print(f"type:    {type(result)}")
\end{verbatim}

Possible output:

\begin{verbatim}
2 ** -2: 0.25
type:    <class 'float'>
\end{verbatim}

This is read as:

\begin{verbatim}
2 ** (-2)
\end{verbatim}

For this chapter, the practical warning is enough:

\begin{quote}
When exponentiation and unary signs appear together, use parentheses unless the intended grouping is visually obvious.
\end{quote}

\subsubsection{Parentheses as Explicit Precedence Control}

Operator precedence decides how an expression is grouped when parentheses are absent.

A simplified numerical precedence ladder is:

\begin{verbatim}
highest
    **                 exponentiation
    +x, -x             unary plus, unary minus
    *, /, //, %         multiplication and division family
    +, -               addition and subtraction
lowest
\end{verbatim}

This ladder is simplified for the arithmetic operators used here. It is enough for the present section.

Examples:

\begin{verbatim}
result1 = 2 + 3 * 4
result2 = (2 + 3) * 4

print(f"2 + 3 * 4:     {result1}")
print(f"(2 + 3) * 4:   {result2}")
\end{verbatim}

Possible output:

\begin{verbatim}
2 + 3 * 4:     14
(2 + 3) * 4:   20
\end{verbatim}

Without parentheses, multiplication happens before addition.

Another example:

\begin{verbatim}
result1 = 20 - 6 // 2
result2 = (20 - 6) // 2

print(f"20 - 6 // 2:     {result1}")
print(f"(20 - 6) // 2:   {result2}")
\end{verbatim}

Possible output:

\begin{verbatim}
20 - 6 // 2:     17
(20 - 6) // 2:   7
\end{verbatim}

Parentheses do not create a new numeric type. They only control expression grouping.

\begin{verbatim}
a = 2 + 3 * 4
b = (2 + 3) * 4

print(f"a: {a}, type={type(a)}")
print(f"b: {b}, type={type(b)}")
\end{verbatim}

Possible output:

\begin{verbatim}
a: 14, type=<class 'int'>
b: 20, type=<class 'int'>
\end{verbatim}

For teaching and production code, parentheses are often better than relying on the reader to remember every precedence rule.

Poor readability:

\begin{verbatim}
score = base + bonus * level - penalty // 2 ** round_count
\end{verbatim}

Clearer:

\begin{verbatim}
score = base + (bonus * level) - (penalty // (2 ** round_count))
\end{verbatim}

A complete small script:

\begin{verbatim}
base = 42
bonus = 5
level = 3
penalty = 16
round_count = 2

score = base + (bonus * level) - (penalty // (2 ** round_count))

print(f"base:        {base}")
print(f"bonus:       {bonus}")
print(f"level:       {level}")
print(f"penalty:     {penalty}")
print(f"round_count: {round_count}")

print()
print(f"score:       {score}")
print(f"type(score): {type(score)}")
\end{verbatim}

Possible output:

\begin{verbatim}
base:        42
bonus:       5
level:       3
penalty:     16
round_count: 2

score:       53
type(score): <class 'int'>
\end{verbatim}

The practical summary is:

\begin{quote}
Numerical operators create result objects through type-defined behavior. Operator precedence controls how expressions are grouped before those operations execute. Parentheses are the clearest way to state the grouping that the programmer intends.
\end{quote}